% AI-written, AI cross-reviewed; no human sign-off; not a claim that RH is proved
% Atlas unit D155. Local manuscript v3, item G.165, heading on printed PDF p. 198.
% FubiniStudyChowHeightFormula₂.fubiniStudy_chowHeight_formula
% paper2/decomposition_tables.tex:1409–1414. Verbatim excerpt; SHA-256 25cc6503a6e3d3ba06ac32475cf8c2d534482565a088b239abc64eae9a338f64.
% Original TeX uses the macros and cross-reference keys of the source manuscript.


% Start of atlas edition annotation. The following block is outside the original excerpt.
\paragraph{Status of named inputs in this reading.}
\noindent\nolinkurl{GlobalTheory.CovTop.PublishedHodgeBundleTheorems.p4} --- record field (\path{Development/AutomorphicArithmetic/ThetaCorrespondence/ThetaPairLifts/HodgeBundleWeights/ToKnownTheoremsOf.lean:11}); not compared.\par
% End of atlas edition annotation.

\begin{formalstatement}[D155]\label{decomp:D155}
The published Hodge-bundle package gives the Fubini–Study expression for the height of the common Chow form.
\end{formalstatement}
{\small\raggedright\noindent\textsf{\bfseries Related reference}\quad \nolinkurl{FubiniStudyChowHeightFormula₂.fubiniStudy_chowHeight_formula}\par\noindent \nolinkurl{Development/AutomorphicArithmetic/ThetaCorrespondence/ThetaPairLifts/HodgeBundleWeights/FubiniStudyChowHeightFormula.lean:17--18}\par\noindent SHA-256 \nolinkurl{0bb29f3202fc1b662285ddf5c72e98141b12a0519d384554a999f7c07a99013b}\par\noindent Dominated 0\quad Disjoint 0\quad terminal.\par}
The Chow form, scalar norm and Fubini–Study normalization refer to specific project constructions.
\hyperref[lit:S036]{Published sources and search S036}.
